\section{Conclusion and Future Work}
\label{sec:conclusion}

This paper presented \mra{}, a configurable and reproducible benchmark and
evaluation framework for onboard autonomous underwater gate racing. Its central
contribution is a controller-agnostic participant interface, based on a common
participant-level information boundary and body-frame action interface, combined
with configurable race management and an independent referee. Controllers use
participant-accessible onboard information, while privileged simulator state is
reserved for ordered-crossing validation, penalties and official race results.

The experiments show that this interface can support heterogeneous circuits,
environmental disturbance, multi-vehicle evaluation and different controller
families. The reference controllers expose condition-dependent behaviour across
the official circuits and under the medium-current condition, while the fleet
experiments demonstrate team-level evaluation and coordination. The recurrent
PPO controller further shows that a learned policy can be integrated through the
same participant-level information boundary, body-frame action interface and
independent referee. Together, these results show that the benchmark is not tied
to a particular controller architecture.

The current evidence remains simulation-based and is limited by simulator
fidelity, a single evaluated medium-current condition, simplified acoustic
models and small fleet configurations. The PPO controller is intentionally a
baseline learning demonstration and can be extended with more advanced learning
algorithms and training strategies. Future work will broaden the disturbance and
fleet evaluations, investigate stronger learning-based controllers and transfer
the benchmark interface to physical underwater experiments.
